\chapter{Yüklemler ve Altküme Operatörleri}
\label{ch:predicates_operators}

%% ============================================================
\section{Konu}
%% ============================================================

\subsection{Altküme Yüklemleri}

$(X, \tau)$ bir topolojik uzay ve $A \subseteq X$ olsun.

\begin{table}[h]
\centering
\begin{tabular}{ll}
\toprule
\textbf{Yüklem} & \textbf{Tanım} \\
\midrule
Açık (open) & $A \in \tau$ \\
Kapalı (closed) & $X \setminus A \in \tau$ \\
Clopen & $A \in \tau$ ve $X \setminus A \in \tau$ \\
Yoğun (dense) & Her boş olmayan $U \in \tau$ için $U \cap A \neq \emptyset$ \\
Hiçbiryerde-yoğun & $\mathrm{int}(\overline{A}) = \emptyset$ \\
\bottomrule
\end{tabular}
\caption{Altküme yüklemleri}
\end{table}

\subsection{Altküme Operatörleri}

\begin{table}[h]
\centering
\begin{tabular}{lll}
\toprule
\textbf{Operatör} & \textbf{Gösterim} & \textbf{Tanım} \\
\midrule
Kapanış & $\overline{A} = \mathrm{cl}(A)$ & $A$'yı içeren en küçük kapalı küme \\
İç & $A^\circ = \mathrm{int}(A)$ & $A$'ya dahil en büyük açık küme \\
Sınır & $\partial A = \mathrm{bd}(A)$ & $\overline{A} \setminus A^\circ$ \\
Türev kümesi & $A' = d(A)$ & $A$'nın birikme noktaları kümesi \\
\bottomrule
\end{tabular}
\caption{Altküme operatörleri}
\end{table}

\noindent\textbf{Birikme noktası:} $x \in A'$ $\iff$ her $U \ni x$ açık için
$U \cap (A \setminus \{x\}) \neq \emptyset$.

\subsection{Komşuluk Sistemi}

\begin{definition}[Komşuluk Sistemi]
$x \in X$ için:
\[
  N(x) = \{N \subseteq X : \exists U \in \tau,\, x \in U \subseteq N\}
\]
\end{definition}

\noindent\textbf{Komşuluk Aksiyomları (N1--N4):}
\begin{enumerate}
  \item[(N1)] $x \in N$, her $N \in N(x)$ için
  \item[(N2)] $X \in N(x)$
  \item[(N3)] $N_1, N_2 \in N(x) \Rightarrow N_1 \cap N_2 \in N(x)$
  \item[(N4)] $N \in N(x)$, $N \subseteq M \Rightarrow M \in N(x)$
\end{enumerate}

%% ============================================================
\section{Teoremler}
%% ============================================================

\begin{theorem}[Kuratowski Kapanış Aksiyomları]
\label{thm:kuratowski}
$\mathrm{cl}: \mathcal{P}(X) \to \mathcal{P}(X)$ operatörü şu dört özelliği sağlar:
\begin{enumerate}
  \item[(K1)] $\mathrm{cl}(\emptyset) = \emptyset$
  \item[(K2)] $A \subseteq \mathrm{cl}(A)$
  \item[(K3)] $\mathrm{cl}(\mathrm{cl}(A)) = \mathrm{cl}(A)$
  \item[(K4)] $\mathrm{cl}(A \cup B) = \mathrm{cl}(A) \cup \mathrm{cl}(B)$
\end{enumerate}
Bu dört özelliği sağlayan her operatör bir topoloji tanımlar.
\end{theorem}

\begin{theorem}[İç-Kapanış Dualitesi]
\label{thm:interior_closure_duality}
Her $A \subseteq X$ için:
\[
  A^\circ = X \setminus \mathrm{cl}(X \setminus A), \qquad
  \mathrm{cl}(A) = X \setminus \mathrm{int}(X \setminus A).
\]
\end{theorem}

\begin{theorem}[Komşuluk Sistemi Round-Trip]
\label{thm:neighborhood_round_trip}
N1--N4'ü sağlayan her $\{N(x)\}_{x \in X}$ ailesi,
$\tau = \{U : \forall x \in U, U \in N(x)\}$ şeklinde tek bir topolojiyi geri üretir.
\end{theorem}

%% ============================================================
\section{Algoritmalar}
%% ============================================================

\begin{algorithm}
\caption{Sonlu Kapanış Hesabı}
\begin{algorithmic}[1]
\Require $A \subseteq X$, $\tau$
\Ensure $\mathrm{cl}(A)$
\State $\textit{closed} \leftarrow \{X \setminus U : U \in \tau\}$
\State $\textit{result} \leftarrow X$
\For{each $C \in \textit{closed}$}
    \If{$A \subseteq C$ \textbf{and} $|C| < |\textit{result}|$}
        \State $\textit{result} \leftarrow C$
    \EndIf
\EndFor
\Return $\textit{result}$
\end{algorithmic}
\end{algorithm}

\noindent\textbf{Karmaşıklık:} $O(|\tau| \cdot |X|)$.

\begin{algorithm}
\caption{Türev Kümesi Hesabı}
\begin{algorithmic}[1]
\Require $A \subseteq X$, $\tau$
\Ensure $d(A)$
\State $\textit{acc} \leftarrow \emptyset$
\For{each $x \in X$}
    \If{her açık $U \ni x$ için $U \cap (A \setminus \{x\}) \neq \emptyset$}
        \State $\textit{acc}.\text{add}(x)$
    \EndIf
\EndFor
\Return $\textit{acc}$
\end{algorithmic}
\end{algorithm}

\noindent\textbf{Karmaşıklık:} $O(|X| \cdot |\tau|)$.

%% ============================================================
\section{pytop API}
%% ============================================================

\begin{lstlisting}[language=Python, caption={Subset operatörleri importları}]
from pytop import (
    analyze_predicate,
    closure_of_subset, interior_of_subset, boundary_of_subset,
    derived_set_of_subset,
    is_open_subset, is_closed_subset, is_dense_subset, is_nowhere_dense_subset,
    neighborhood_system, character_at_point, analyze_neighborhood_system,
)
\end{lstlisting}

\begin{table}[h]
\centering
\begin{tabular}{lll}
\toprule
\textbf{Fonksiyon} & \textbf{İmza} & \textbf{Açıklama} \\
\midrule
\texttt{analyze\_predicate} & \texttt{(space, predicate, subset)} & Yüklem sorgula \\
\texttt{closure\_of\_subset} & \texttt{(space, subset)} & Kapanış \\
\texttt{interior\_of\_subset} & \texttt{(space, subset)} & İç \\
\texttt{boundary\_of\_subset} & \texttt{(space, subset)} & Sınır \\
\texttt{derived\_set\_of\_subset} & \texttt{(space, subset)} & Türev kümesi \\
\texttt{neighborhood\_system} & \texttt{(carrier, topology, point)} & $N(x)$ \\
\texttt{character\_at\_point} & \texttt{(carrier, topology, point)} & $\chi(x)$ \\
\texttt{analyze\_neighborhood\_system} & \texttt{(carrier, topology, point=None)} & Tam analiz \\
\bottomrule
\end{tabular}
\caption{Altküme operatör fonksiyonları}
\end{table}

\noindent\textbf{Not:} \texttt{neighborhood\_system} ve ilgili fonksiyonlar \texttt{carrier} (liste)
ve \texttt{topology} (liste) alır. Space nesnesi için: \texttt{list(space.carrier)},
\texttt{list(space.topology)}.

%% ============================================================
\section{Örnekler}
%% ============================================================

\subsection{Örnek 1: Kapanış ve İç — Sierpiński Uzayı}

\begin{lstlisting}[language=Python]
from pytop import sierpinski_space, closure_of_subset, interior_of_subset, boundary_of_subset

s = sierpinski_space()
print("cl({0})  =", closure_of_subset(s, {0}).value)
print("cl({1})  =", closure_of_subset(s, {1}).value)
print("int({0}) =", interior_of_subset(s, {0}).value)
print("int({1}) =", interior_of_subset(s, {1}).value)
print("bd({1})  =", boundary_of_subset(s, {1}).value)
\end{lstlisting}

\begin{lstlisting}[caption={Çıktı}]
cl({0})  = frozenset({0})
cl({1})  = frozenset({0, 1})
int({0}) = frozenset()
int({1}) = frozenset({1})
bd({1})  = frozenset({0})
\end{lstlisting}

$\{0\}$'ın kapanışı kendisidir: $\{0\} = X \setminus \{1\}$ zaten kapalıdır.
$\{1\}$'in kapanışı $X$'tir: $\{1\}$ kapalı değildir.

\subsection{Örnek 2: Türev Kümesi}

\begin{lstlisting}[language=Python]
from pytop import sierpinski_space, derived_set_of_subset, is_nowhere_dense_subset

s = sierpinski_space()
print("d({0}) =", derived_set_of_subset(s, {0}).value)
print("d({1}) =", derived_set_of_subset(s, {1}).value)
print("{0} nowhere dense? =", is_nowhere_dense_subset(s, {0}).status)
\end{lstlisting}

\begin{lstlisting}[caption={Çıktı}]
d({0}) = frozenset()
d({1}) = frozenset({0})
{0} nowhere dense? = true
\end{lstlisting}

$d(\{1\}) = \{0\}$: $0$'ı içeren her $\{0,1\}$, $\{1\} \setminus \{0\} = \{1\}$ ile kesişir.

\subsection{Örnek 3: Komşuluk Sistemi}

\begin{lstlisting}[language=Python]
from pytop import sierpinski_space, neighborhood_system, character_at_point

s = sierpinski_space()
carrier = list(s.carrier)
topology = list(s.topology)
print("N(0) =", neighborhood_system(carrier, topology, 0).value)
print("N(1) =", neighborhood_system(carrier, topology, 1).value)
print("chi(0) =", character_at_point(carrier, topology, 0).value)
print("chi(1) =", character_at_point(carrier, topology, 1).value)
\end{lstlisting}

\begin{lstlisting}[caption={Çıktı}]
N(0) = [['0', '1']]
N(1) = [['1'], ['0', '1']]
chi(0) = 1
chi(1) = 2
\end{lstlisting}

$N(0) = \{X\}$: $0$ yalnızca $X$'te açık. $N(1) = \{\{1\}, X\}$: yerel baz büyüklükleri
$\chi(0)=1$, $\chi(1)=2$.

%% ============================================================
\section{Alıştırmalar}
%% ============================================================

\subsection*{Kodlama Alıştırmaları}
\begin{enumerate}[label=\textbf{K\arabic*.}]
  \item İndirgenmiş iki-noktalı topolojide her noktanın komşuluk sistemini hesaplayın.
  \item $X = \{1,2,3,4\}$, $\tau = \{\emptyset, \{1,2\}, \{3,4\}, X\}$ topolojisinde
        $\mathrm{bd}(\{1,2,3\})$ ve $\mathrm{cl}(\{1,2,3\})$ hesaplayın.
  \item \texttt{finite\_chain\_space(4)} zincirinde $\{1,2\}$'nin iç, kapanış ve sınır
        kümelerini bulun.
\end{enumerate}

\subsection*{Teori Alıştırmaları}
\begin{enumerate}[label=\textbf{T\arabic*.}]
  \item $A^\circ = X \setminus \mathrm{cl}(X \setminus A)$ eşitliğini Kuratowski
        aksiyomlarını kullanarak kanıtlayın.
  \item $A$'nın yoğun olması için gerek ve yeter koşulun $\mathrm{cl}(A) = X$ olduğunu
        gösterin.
\end{enumerate}
